\section{Scope and intended use}
\label{sec:scope}

\subsection{What the tool is}

\Tool{} is a screening-level, multi-criteria decision-support instrument. It
converts a structured description of an urban site and one or more proposed
nature-based interventions into comparable scores and indicative quantitative
estimates, so that early planning and investment decisions can be made on a
transparent, evidence-referenced basis.

A user describes a site through six input groups, selects one or more of the
\num{121} catalogue entries the tool offers for that site
(\Cref{sec:typologies}), and receives a Heat Priority Index, a Cooling
Potential Score with an indicative daytime pedestrian-level air temperature
reduction range, derived cooling-energy and greenhouse-gas estimates, cost
outputs where cost data has been supplied, equity and co-benefit scores, a
single NbS Cooling Opportunity Score for comparing options, per-block
confidence ratings, and an itemised list of every default the engine applied.
Where several interventions are selected, each is scored and reported
individually and the package's headline outputs are combined under the rules of
\Cref{sec:packages}, which never sum a temperature. Interventions and packages
for the same site can be assessed and compared side by side.

\subsection{What the tool is not}
\label{sec:not}

The following exclusions are definitional, not caveats to be read past.

\begin{itemize}
  \item \textbf{It is not a microclimate simulation.} It does not solve an
  energy balance, does not represent site geometry, and does not predict the
  temperature at a specific location at a specific time. Its outputs are not
  comparable to ENVI-met-class or CFD results and must never be presented as
  such.

  \item \textbf{It is not a building energy model.} Cooling-energy savings are
  derived from an estimated air temperature reduction through a published
  temperature--electricity-demand sensitivity (\Cref{sec:energy}). No building
  is modelled.

  \item \textbf{It is not a planting design or species selection tool.} The
  evidence that species and leaf area index materially change delivered cooling
  \parencite{rahman2020} is used only to justify that canopy condition
  modulates performance. The tool gives no species guidance.

  \item \textbf{It is not a regulatory, engineering, or approval instrument.}
  Outputs are indicative ranges for comparison and prioritisation. They are not
  design specifications and they are not contractual performance guarantees.

  \item \textbf{It does not replace site investigation or local expertise.} It
  structures a judgment; it does not substitute for one.
\end{itemize}

\subsection{Intended users and decisions}

Three user groups are supported: municipal planners and public agencies;
developers and consultants; and researchers and non-governmental
organisations. The tool is designed to inform three questions:

\begin{enumerate}
  \item Does this site merit attention at all, given how hot it is and who is
  exposed?
  \item Which of several candidate interventions best suits this particular
  site?
  \item How does a proposed intervention perform simultaneously across cooling,
  energy, climate, cost, equity, and co-benefit dimensions?
\end{enumerate}

Each of these is a comparative question. The tool is built to rank and to
compare, and its outputs should be interpreted comparatively. An Opportunity
Score of 72 is meaningful principally in relation to the score another option
or another site receives under the same methodology version; it is not a
physical measurement.

\subsection{The five design commitments}
\label{sec:commitments}

Five commitments constrain the methodology. They are listed here because
several of them cost the tool apparent capability, and a reader should
understand that those costs were chosen.

\begin{table}[H]
\centering
\caption{Design commitments and their consequences for the methodology.}
\label{tab:commitments}
\small
\begin{tabularx}{\textwidth}{@{}l>{\raggedright\arraybackslash}X@{}}
\toprule
\textbf{Commitment} & \textbf{Consequence} \\
\midrule
Transparency &
Every score has exactly one documented formula, reproduced in this paper. No
machine learning, no opaque calibration, no undisclosed adjustment. A user who
disagrees with a number can locate the rule that produced it. \\
\addlinespace
Determinism &
Identical inputs under an identical methodology version always yield identical
outputs. There are no stochastic components anywhere in the scoring or
reporting path, including the recommendation text, which is composed from
templates rather than generated. \\
\addlinespace
Evidence traceability &
Every performance value cites a source. Values that could not be grounded in
retrieved literature are either omitted --- as unit costs are
(\Cref{sec:costs}) --- or shipped with an explicit low-confidence flag. \\
\addlinespace
Honest uncertainty &
All quantitative outputs are ranges. Every result carries a confidence rating,
computed separately per output block. Missing data reduces confidence rather
than being silently imputed, and every applied default is itemised in the
result. \\
\addlinespace
Auditability &
The methodology lives in versioned configuration files, not in code. Every
result records the methodology version that produced it, and the evidence
rules are enforced by continuous integration rather than by review convention
(\Cref{sec:implementation}). \\
\bottomrule
\end{tabularx}
\end{table}

The commitment to honest uncertainty is the one that most visibly reduces the
tool's apparent authority. A screening instrument that reported single numbers
would look more decisive and would be easier to put in a slide. The narrowest
adopted cooling envelope spans a factor of three; the widest with a non-zero
floor, the large open water body, spans a factor of thirty because remote
sensing and field measurement disagree by an order of magnitude
(\Cref{sec:water-body}); and four archetypes carry a genuine zero floor,
for which no ratio is defined at all. The climate-differentiated
estimates
in \textcite{keravec2026} carry uncertainty intervals that in several cases
approach or exceed $\pm\,$\qty{1}{\degreeCelsius}, and in the case of street
trees in arid climates reach $\pm\,$\qty{1.4}{\degreeCelsius} on a central
estimate of \qty{1.7}{\degreeCelsius}; a point estimate would imply a
precision the evidence does not support. The ranges are the honest
representation of what is known.
